% =====================================================================
% FaaS-MAPG: sequential better-response dynamics of an exact potential
% game, with an epsilon-Nash certificate.
% Meant to be \input{} (or pasted) into the paper. Assumes the paper's
% notation (N, F, C_i^f, rho_i, omega_i^f, beta_{ij}^f, gamma_i^f, ...).
% Cross-references are plain text; convert to \ref{} in the host paper.
% =====================================================================

\section{FaaS-MAPG: multi-agent potential-game better response}
\label{sec:mapg}

\textsc{FaaS-MAPG} (Multi-Agent Potential Game) revisits the same
Gauss-Seidel sweep skeleton as \textsc{FaaS-MABR}, but replaces the
score-greedy, commit-always response with a response rule chosen so that the
sum of node utilities is an \emph{exact potential function}: every accepted
unilateral move raises the sum by exactly the mover's utility gain. This is
qualitatively different from the two other coordination families in this
paper. \textsc{FaaS-MADeA} and \textsc{FaaS-MALD} coordinate through prices
(an $\varepsilon$-auction and a Lagrangian dual, respectively) and their
certificates bound a fixed-capacity inner problem, not the joint outer game.
\textsc{FaaS-MABR} is a greedy heuristic: nothing prevents a sweep from
lowering the objective, and it offers no equilibrium notion. \textsc{FaaS-MAPG}
instead accepts a move only when it strictly improves the mover's own utility
by more than a tolerance $\varepsilon>0$, which (Proposition~\ref{prop:mapg-potential})
makes the welfare function $\Phi=\sum_i u_i$ increase by exactly the same
amount, is monotonically non-decreasing, and yields a finite-time certificate
(Theorem~\ref{thm:mapg-termination}) that the terminal state is an
$\varepsilon$-Nash equilibrium with respect to the implemented move class.

% =====================================================================
% BEGIN self-contained notation recap.
% Reproduces the relevant rows of the paper's Table 2 (FRALB) and Table 3
% (auction) plus Eqs. (1),(2),(12),(13), so the note is readable on its own.
% REMOVE this whole subsection when inserting into the paper, where the
% notation and these equations are already defined.
% =====================================================================
\subsection{Notation and capacity model}
\label{sec:mapg-notation}

Each control period $t$ runs coordination iterations (\emph{sweeps})
$h=0,1,\dots,H_{\max}$. Before coordination, every node $i$ solves a local
problem~(P2) that, for each function $f\in\mathcal{F}_i$, splits the incoming
workload $\lambda_i^f$ into locally processed load $x_i^f$, horizontally
offloaded demand $\omega_i^f$, and Cloud-forwarded load $z_i^f$, and fixes the
number of replicas $r_i^f$. Replicas obey the memory budget and provide
capacity
\begin{equation}
  \sum_{f\in\mathcal{F}_i} r_i^f\,\mathrm{RAM}^f \le \mathrm{RAM}_i,
  \qquad
  \mathrm{Cap}_i^f = \frac{r_i^f\,U^f_{\max}}{D_i^f}\ \text{[req/s]} ,
  \tag{1,2}
\end{equation}
from which each node advertises, on a neighbour-accessible blackboard, its
residual execution capacity and its memory slack:
\begin{equation}
  C_i^f(h) = \max\!\big\{0,\ \mathrm{Cap}_i^f(h) - x_i^f(h)\big\},
  \qquad
  \rho_i(h) = \mathrm{RAM}_i - \sum_{f\in\mathcal{F}_i} r_i^f(h)\,\mathrm{RAM}^f \ge 0 .
  \tag{12,13}
\end{equation}
Table~\ref{tab:mapg-notation} collects the symbols used below.

\begin{table}[t]
\centering
\caption{Notation used in this section (subset of the paper's Tables~2--3).}
\label{tab:mapg-notation}
\small
\begin{tabular}{@{}l l p{0.6\linewidth}@{}}
\toprule
\multicolumn{3}{@{}l}{\textit{Sets and players}}\\
Nodes / neighbours & $\mathcal{N}$, $N_i$ & computing nodes, treated as players; one-hop neighbours of $i$.\\
Functions & $\mathcal{F}_i$ & functions deployable on $i$; $\mathcal{F}=\bigcup_i \mathcal{F}_i$.\\
Move order & $\sigma$ & the sweep's visiting order over $\mathcal{N}$ (fixed or a per-sweep random permutation).\\
\midrule
\multicolumn{3}{@{}l}{\textit{Parameters and weights}}\\
Input workload & $\lambda_i^f$ & incoming request rate for $f$ at $i$ [req/s]; also the utility normalizer.\\
Memory cap / demand & $\mathrm{RAM}_i$, $\mathrm{RAM}^f$ & node memory; per-replica memory of $f$ [MB].\\
Service demand & $D_i^f$ & average execution time of $f$ at $i$ [s].\\
Max utilization & $U^f_{\max}$ & stability threshold, $U^f_{\max}\in[0,1)$.\\
Local / offload / Cloud profit & $\alpha_i^f$, $\beta_{ij}^f$, $\gamma_i^f$ & profit of local exec.; profit of offloading $i\!\to\! j$; Cloud-offloading penalty (per req/s).\\
Acceptance tolerance & $\varepsilon>0$ & minimum utility gain required to commit a move.\\
Max iterations / time limit & $H_{\max}$, $T_{\max}$ & sweep and wall-clock stop guards.\\
\midrule
\multicolumn{3}{@{}l}{\textit{Strategy and state at sweep $h$}}\\
Node strategy & $s_i=(r_i,x_i,\omega_i,y_{i\cdot})$ & replicas, local load, offloading target, and outbound routing row.\\
Replicas & $r_i^f$ & number of $f$ replicas on $i$, $r_i^f\in\mathbb{N}_0$.\\
Local / forwarded / Cloud load & $x_i^f$, $y_{ij}^f$, $z_i^f$ & locally served; forwarded $i\!\to\! j$; Cloud-offloaded [req/s].\\
Committed inbound & $\bar y_{mi}^f$ & the current $y_{mi}^f$ of every node $m\ne i$, frozen while $i$ moves.\\
Accessible capacity cap & $\omega_i^{\mathrm{ub},f}$ & the residual capacity $i$ may claim from its neighbours in this move.\\
Capacity / residual / slack & $\mathrm{Cap}_i^f$, $C_i^f$, $\rho_i$ & exec.\ capacity (Eq.~2); residual exec.\ capacity (Eq.~12); memory slack (Eq.~13).\\
Node utility & $u_i$ & node $i$'s share of the centralized FRALB objective (Eq.~\eqref{eq:mapg-utility}).\\
Potential / welfare & $\Phi$ & $\Phi=\sum_i u_i$, the centralized objective; the exact potential (Prop.~\ref{prop:mapg-potential}).\\
\bottomrule
\end{tabular}
\end{table}
% ===================== END self-contained notation recap =====================

\subsection{Game formulation}
\label{sec:mapg-game}

Players are the nodes $\mathcal{N}$. Node $i$'s strategy is
$s_i=(r_i,x_i,\omega_i,y_{i\cdot})$: its replica counts, locally served load,
offloading target, and the row $y_{i\cdot}^f=(y_{ij}^f)_{j\in N_i}$ of load it
routes to neighbours (Cloud-forwarded load $z_i^f$ is not a free coordinate;
it is the flow-conservation residual $z_i^f=\max\{0,\ell_i^f-x_i^f-\sum_j y_{ij}^f\}$
of $i$'s own served load $\ell_i^f$). Node $i$'s utility is its share of the
centralized FRALB objective, i.e. the terms of the objective that depend on
$i$'s own coordinates:
\begin{equation}
  u_i(s) = \sum_{f\in\mathcal{F}_i}
    \frac{\alpha_i^f x_i^f + \sum_{j\in N_i}\beta_{ij}^f y_{ij}^f - \gamma_i^f z_i^f}
         {\lambda_i^f}.
  \label{eq:mapg-utility}
\end{equation}
Summing $u_i$ over all nodes recovers exactly the centralized objective
$\Phi(s)=\sum_i u_i(s)$: each $(\alpha,\beta,\gamma)$-weighted term is counted
once, against the node that owns the corresponding row of $x$, $y$, or $z$.

\paragraph{Unilateral-move rules.} A move by node $i$ changes only $i$'s
strategy $s_i$; every other node's strategy $s_{-i}$, including its outbound
row $y_{m\cdot}$ for $m\ne i$, is held fixed. Two rules make this well
defined and keep $\Phi$ an exact potential:
\begin{enumerate}
\item \textbf{Inbound commitments are kept.} Load already routed to $i$ by
  other nodes, $\bar y_{mi}^f$ for $m\ne i$, must still be served by $i$ after
  the move; $i$'s local problem enforces $D_i^f\big(x_i^f+\sum_m \bar y_{mi}^f\big)\le r_i^f U_{\max}^f$
  (utilization constraints \texttt{utilization\_equilibrium\_pg} /
  \texttt{\_2\_pg} of \texttt{LSP\_pg} in \texttt{models/sp.py}), so a move
  can never strand another node's already-placed load.
\item \textbf{Only advertised residual capacity is claimable, and $i$'s own
  row is released first.} Before computing $i$'s move, $i$'s current outbound
  row $y_{i\cdot}$ is zeroed out of the residual-capacity ledger (so $i$ does
  not compete with the capacity it is about to give up), and $i$ is capped to
  the resulting residual capacity of eligible neighbours,
  $\omega_i^{\mathrm{ub},f}=\sum_{j\in N_i:\ \beta_{ij}^f>-\gamma_i^f}\tilde C_j^f$,
  enforced by the \texttt{cap\_offloading} constraint $\omega_i^f\le\omega_i^{\mathrm{ub},f}$.
\item \textbf{No ping-pong (FRALB feasibility).} FRALB forbids a node from
  simultaneously sending and receiving the same function. Accordingly, a
  neighbour $j$ currently offloading $f$ (i.e., $\sum_k y_{jk}^f>0$) is
  excluded from the seller set for $f$, and the mover sets
  $\omega_i^{\mathrm{ub},f}=0$ for every $f$ it holds inbound commitments
  for ($\sum_m \bar y_{mi}^f>0$). Since $y=0$ initially, these two rules
  inductively keep every intermediate state ping-pong-free. Both only
  shrink the feasible move set, so the exact-potential argument below is
  unaffected.
\end{enumerate}
Because $i$'s move only ever reallocates $i$'s own row $y_{i\cdot}$ within
capacity every other node has already advertised as unused, and every other
node's coordinates $s_{-i}$ (including $\bar y_{mi}$, which $i$ must
continue to serve) are literally unchanged by the move, no other node's
utility term changes: $u_m(s_i',s_{-i})=u_m(s_i,s_{-i})$ for all $m\ne i$.

\subsection{Exact potential}
\label{sec:mapg-potential}

\begin{proposition}[FaaS-MAPG is an exact potential game]
\label{prop:mapg-potential}
Under the unilateral-move rules of \S\ref{sec:mapg-game}, the welfare function
$\Phi(s)=\sum_{k\in\mathcal{N}} u_k(s)$ is an exact potential for the game:
for every node $i$, every strategy profile $s$, and every unilateral
deviation $s_i'$,
\[
  \Phi(s_i',s_{-i}) - \Phi(s_i,s_{-i}) = u_i(s_i',s_{-i}) - u_i(s_i,s_{-i}).
\]
\end{proposition}
\begin{proof}
Write $\Phi(s_i,s_{-i})=u_i(s_i,s_{-i})+\sum_{m\ne i}u_m(s_i,s_{-i})$ and
likewise for $\Phi(s_i',s_{-i})$. Subtracting,
\[
  \Phi(s_i',s_{-i})-\Phi(s_i,s_{-i})
  = \big[u_i(s_i',s_{-i})-u_i(s_i,s_{-i})\big]
  + \sum_{m\ne i}\big[u_m(s_i',s_{-i})-u_m(s_i,s_{-i})\big].
\]
Each $u_m$, $m\ne i$, is a function of $m$'s own coordinates
$(x_m,y_{m\cdot},z_m)$ only (Eq.~\eqref{eq:mapg-utility}). A unilateral move
by $i$ leaves $s_{-i}$ untouched by construction, and rule~1 guarantees that
$m$'s served load $\bar y_{mi}^f$ from $i$'s row, hence $m$'s own $x_m$,
$y_{m\cdot}$, $z_m$, do not have to and do not change to remain feasible
(the move only reduces or reallocates what $i$ sends, which $m$ does not
own). Consequently $u_m(s_i',s_{-i})=u_m(s_i,s_{-i})$ for every $m\ne i$, the
sum of differences vanishes, and
$\Phi(s_i',s_{-i})-\Phi(s_i,s_{-i})=u_i(s_i',s_{-i})-u_i(s_i,s_{-i})$.
\end{proof}

This is the FaaS instance of the standard construction of an exact potential
game from an additively separable global objective in which each player's
payoff is exactly its own contribution to that objective
\citep{monderershapley1996}; the resulting $\Phi$ coincides with the FRALB
welfare, and the game is, structurally, a congestion-style game over shared
capacity resources in the sense of \citet{rosenthal1973}, restricted here to
the implemented move class rather than unrestricted resource choice.

\subsection{Move rule: proposal MILP, exact split, and $\varepsilon$-acceptance}
\label{sec:mapg-move}

A single node move (function \texttt{node\_move} in
\texttt{decentralized\_potentialgame.py}) has three stages.

\paragraph{1. Proposal MILP.} Node $i$ re-solves its local problem
(\texttt{LSP\_pg}, or \texttt{LSP\_pg\_fixedr} when replicas are pinned to a
reference plan) with committed inbound flows $\bar y_{mi}^f$ fixed as a
parameter and the offloading target capped at $\omega_i^{\mathrm{ub},f}$
(\S\ref{sec:mapg-game}). This MILP is \emph{P2 with committed inbound and a
capacity cap}: it prices the aggregate offload $\omega_i^f$ using the same
weighted objective as~(P2), i.e.\ it treats every unit of $\omega_i^f$ as
interchangeable rather than pricing each destination pair $\beta_{ij}^f$
individually. Its output is a proposed $x_i$, $r_i$, and offload amount
$\omega_i$; it does \emph{not} decide the destination split.

\paragraph{2. Fractional-knapsack split by descending $\beta$.} The proposed
$\omega_i^f$ is placed into a routing row by \texttt{split\_omega}: eligible
neighbours $j$ (those with $\beta_{ij}^f>-\gamma_i^f$ and unused ledger
capacity) are ranked by descending true pair-specific profit $\beta_{ij}^f$,
ties broken by ascending node index $j$, and filled greedily up to each
neighbour's remaining ledger capacity until $\omega_i^f$ is exhausted or
sellers run out; any residual is left to the Cloud ($z_i^f$). Because the
per-unit payoff $\beta_{ij}^f$ is fixed and the feasible region is a simple
capacity knapsack (a matroid / linear-program vertex problem), this greedy
split is the exact utility-maximizing response to the proposed aggregate
$\omega_i^f$: no other split of the same total achieves a strictly higher
$\sum_j\beta_{ij}^fy_{ij}^f$ term of $u_i$.

\paragraph{3. $\varepsilon$-acceptance on the true
utility.} Only after the split is constructed is the candidate utility $u_i$
recomputed (\texttt{compute\_node\_utility}, using the true per-pair
$\beta_{ij}^f$, unlike the proposal MILP's aggregate pricing) and compared to
the utility of $i$'s current strategy. The move is committed to $(x,y,r)$
iff $\Delta u_i = u_i^{\mathrm{new}}-u_i^{\mathrm{old}} > \varepsilon$;
otherwise $i$'s strategy is left unchanged for this sweep. This is why the
proposal MILP's aggregate pricing does not compromise Proposition~\ref{prop:mapg-potential}:
acceptance is gated on the true, exactly-decomposed $u_i$, not on the
proposal's own objective value.

\begin{algorithm}[H]
\caption{\textsc{FaaS-MAPG} --- one better-response sweep}
\label{alg:mapg-sweep}
\begin{algorithmic}[1]
\Require state $(x,y,r)$, ledger seed $C(h)$, order $\sigma$, tolerance $\varepsilon$
\State $n_{\mathrm{accepted}}\gets 0$, $\Delta\Phi\gets 0$, memory bids $\gets\emptyset$
\For{each node $i$ in order $\sigma$}
  \State $u_i^{\mathrm{old}}\gets u_i(x,y,z(x,y))$
  \State release $i$'s own row: $y_{i\cdot}\gets0$ in a working copy; rebuild ledger $\tilde C$ from that copy
  \State $\omega_i^{\mathrm{ub},f}\gets\sum_{j\in N_i:\beta_{ij}^f>-\gamma_i^f}\tilde C_j^f$ for all $f$
  \State solve proposal MILP (\texttt{LSP\_pg}) with $\bar y_{\cdot i}$ fixed and cap $\omega_i^{\mathrm{ub}}$; obtain $x_i',r_i',\omega_i'$
  \State $y_{i\cdot}'\gets\textsc{SplitOmega}(\omega_i',\tilde C)$: descending $\beta_{ij}^f$, ties by ascending $j$
  \State $u_i^{\mathrm{new}}\gets u_i(x',y',z(x',y'))$ with $x_i\gets x_i',\,y_{i\cdot}\gets y_{i\cdot}'$
  \If{$u_i^{\mathrm{new}}-u_i^{\mathrm{old}}>\varepsilon$}
    \State commit $x_i\gets x_i'$, $y_{i\cdot}\gets y_{i\cdot}'$, $r_i\gets r_i'$;\quad $n_{\mathrm{accepted}}{+}{+}$;\quad $\Delta\Phi\mathrel{+}=u_i^{\mathrm{new}}-u_i^{\mathrm{old}}$
  \EndIf
  \State for every $f$ with unplaced appetite $\omega_i'^f-\sum_j y_{ij}'^f>\text{tol}$, emit memory bids to neighbours $j$ with $\rho_j\ge\mathrm{RAM}^f$ and $\beta_{ij}^f>-\gamma_i^f$
\EndFor
\State \Return $n_{\mathrm{accepted}}$, $\Delta\Phi$, memory bids
\end{algorithmic}
\end{algorithm}

The outer loop (\texttt{potential\_game\_sweep} inside the per-timestep loop
of \texttt{decentralized\_potentialgame.py}) repeats the sweep, starts
replicas from any collected memory bids (\S\ref{sec:mapg-memory}), and checks
the stopping condition (\texttt{check\_pg\_stopping}) after each sweep.

\subsection{Finite termination and $\varepsilon$-Nash certificate}
\label{sec:mapg-termination}

\begin{theorem}[Finite termination at an $\varepsilon$-Nash equilibrium]
\label{thm:mapg-termination}
Let $\Phi_0$ be the potential of the initial state and let
$\Phi_{\max}<\infty$ be any finite upper bound on $\Phi$ over the run's
feasible region (finite since $x,y,z,r$ are bounded by finite loads,
capacities, and memory). Under the move rules of \S\ref{sec:mapg-game} and
the acceptance rule of \S\ref{sec:mapg-move}:
\begin{enumerate}
\item the number of \emph{accepted} moves over the whole run is at most
  $(\Phi_{\max}-\Phi_0)/\varepsilon$;
\item the number of replica-expansion events (\S\ref{sec:mapg-memory}) is
  finite, bounded by $\sum_i\lfloor \mathrm{RAM}_i/\min_f\mathrm{RAM}^f\rfloor$
  total replicas system-wide;
\item hence, absent the max-iterations or time-limit guards, the algorithm
  terminates after finitely many sweeps at a sweep with zero accepted moves
  and zero new replicas, and the terminal state $s^\star$ is an
  $\varepsilon$-Nash equilibrium with respect to the implemented move class:
  for every node $i$, the move that \emph{stage 1--2 of \S\ref{sec:mapg-move}
  would propose} from $s^\star$ satisfies $u_i(s_i',s_{-i}^\star)-u_i(s^\star)\le\varepsilon$.
\end{enumerate}
\end{theorem}
\begin{proof}
(1) By Proposition~\ref{prop:mapg-potential}, an accepted move by $i$
satisfies $\Phi(s')-\Phi(s)=u_i(s')-u_i(s)>\varepsilon$, so $\Phi$ strictly
increases by more than $\varepsilon$ at every accepted move and is
non-decreasing at every rejected move (rejected moves leave the state, hence
$\Phi$, unchanged). Since $\Phi$ starts at $\Phi_0$ and is bounded above by
$\Phi_{\max}$, at most $(\Phi_{\max}-\Phi_0)/\varepsilon$ accepted moves can
occur before $\Phi$ would have to exceed $\Phi_{\max}$; the count is hence
finite and bounded by that quantity.
(2) Each replica-expansion event (\texttt{start\_additional\_replicas},
reused unchanged from \textsc{FaaS-MADeA}) consumes strictly positive memory
slack $\rho_j$ at some node $j$ and never removes a replica, so the total
number of replicas placed system-wide is monotonically non-decreasing and
bounded above by the total memory divided by the smallest per-function
memory footprint; expansion events are therefore finite in number. Because
$\Phi$ is the centralized objective and replicas only ever relax capacity
constraints for subsequent moves (they do not change $x,y,z$ of the sweep
that placed them), an expansion event does not decrease $\Phi$.
(3) Both accepted moves and replica expansions are finite in number and each
sweep with neither is a fixed point of \texttt{check\_pg\_stopping}
(condition \texttt{n\_accepted == 0 and n\_new\_replicas == 0}), so, absent
the two runtime guards, some finite sweep reaches this fixed point (each
non-fixed-point sweep strictly advances one of the two finite counters). At
that terminal state $s^\star$, sweep~\S\ref{sec:mapg-move} was run for every
node and rejected the proposed move, i.e.\ $u_i(s_i',s_{-i}^\star)-u_i(s^\star)\le\varepsilon$
for the specific $s_i'$ that stages 1--2 constructed from $s^\star$ for every
$i$; this is the definition of an $\varepsilon$-Nash equilibrium restricted
to deviations reachable by that move class.
\end{proof}

\paragraph{Scope of the certificate.} The certificate in part~3 is honest
about what it does \emph{not} claim: it is not an $\varepsilon$-Nash
equilibrium of the unrestricted game (arbitrary $s_i'$), only with respect to
the specific move class implemented --- the MILP proposal of stage~1 (which
prices aggregate offloading, not per-pair $\beta$) composed with the greedy
split of stage~2. A node might in principle find a strategy outside this
move class (e.g. a split that is not the descending-$\beta$ knapsack applied
to the MILP's own $\omega_i$, or a jointly different $(x_i,\omega_i)$ pair
that the aggregate-priced MILP does not surface) with a strictly larger
utility gain. The greedy split (stage~2) is exact given the proposed
$\omega_i$, so the residual freedom is entirely in whether the proposal MILP
explores the same $(x_i,\omega_i)$ vertex a per-pair-$\beta$-aware
formulation would; this is the same aggregation gap discussed for the
\textsc{FaaS-MADiG}/\textsc{FaaS-MABR} local problem~(P2) elsewhere in the
paper. In practice, when the max-iterations or time-limit guard fires first,
\texttt{check\_pg\_stopping} reports that reason instead and the run carries
no equilibrium certificate for that timestep.

\subsection{Memory market}
\label{sec:mapg-memory}

Unplaced appetite after the fractional-knapsack split triggers the same
memory-bid mechanism as \textsc{FaaS-MADeA}: node $i$ bids to every neighbour
$j$ with unused memory slack $\rho_j\ge\mathrm{RAM}^f$ and
$\beta_{ij}^f>-\gamma_i^f$; \texttt{start\_additional\_replicas} allocates new
replicas at bid-receiving sellers proportionally to demand fraction, subject
to $\rho_j$. Two properties transfer unchanged from \textsc{FaaS-MADeA} and
matter for Theorem~\ref{thm:mapg-termination}: (i) the market is
\emph{memory-bounded} — replicas are only ever added within $\rho_j\ge0$, so
the process cannot run indefinitely; and (ii) the market is
\emph{$\Phi$-invariant} in the sense of never decreasing $\Phi$ — a seller
who opens a new replica does so purely to relax a capacity constraint seen by
some other node's future proposal MILP, and the transaction has no direct
$\beta$-weighted utility term for the seller itself, so a selfish seller has
no incentive to open a replica on its own account and no disincentive
against doing so when bid (the reward, if any, accrues to the buyer's
$\beta_{ij}^f y_{ij}^f$ term in a later sweep, not to the seller's utility
directly). The mechanism is therefore reused as an exogenous capacity-relief
step layered on top of the potential-game dynamics, not itself a move subject
to $\varepsilon$-acceptance.

\subsection{The two variants}
\label{sec:mapg-variants}

\textsc{FaaS-MAPG-S} sweeps nodes in a fixed ascending order (deterministic,
reproducible). \textsc{FaaS-MAPG-R} draws a fresh random node permutation
each sweep from a per-run generator seeded once. Both variants share the same
move rule, acceptance test, and termination guarantee: the order only affects
\emph{which} sequence of moves is realized and hence the specific
$\varepsilon$-Nash equilibrium reached (potential games can have multiple
equilibria), not whether termination and the certificate hold.

\subsection{Stopping conditions}
\label{sec:mapg-stopping}

\texttt{check\_pg\_stopping} reports exactly one of three reasons per
timestep, checked in order: (i) \texttt{"max iterations reached"} if the
sweep count reaches the configured cap; (ii) \texttt{"reached time limit"} if
accumulated wall-clock/solver runtime exceeds the configured limit; (iii)
\texttt{"epsilon-Nash equilibrium certified"} if a sweep accepts zero moves
and starts zero new replicas — the fixed point of
Theorem~\ref{thm:mapg-termination}. Only reason (iii) carries the equilibrium
certificate of \S\ref{sec:mapg-termination}; (i) and (ii) are runtime guards
that may fire before the certificate is reached, in which case the run's
terminal state carries no such guarantee.

\subsection{Positioning with respect to the literature and the paper's other coordination families}
\label{sec:mapg-related}

\textsc{FaaS-MAPG} shares the Gauss-Seidel sequential-sweep skeleton with
\textsc{FaaS-MABR}: both visit nodes in order $\sigma$, re-solve a capped
local MILP, and place offload greedily by descending $\beta_{ij}^f$. The
difference is entirely in the acceptance rule: \textsc{FaaS-MABR} commits the
proposed move unconditionally (score-greedy, commit-always), which offers no
monotonicity guarantee on the objective; \textsc{FaaS-MAPG} gates commitment
on a strict, true-utility improvement, which (Proposition~\ref{prop:mapg-potential},
Theorem~\ref{thm:mapg-termination}) makes $\Phi$ monotone non-decreasing and
yields a finite-time, certified stopping condition that \textsc{FaaS-MABR}
does not have.

Against \textsc{FaaS-MALD}, the certificates answer different questions.
\textsc{FaaS-MALD}'s $LB$/$UB$ pair (see the companion \texttt{faas-mald-note},
\S\S~``Coordination transportation LP''--``Positioning with respect to the
literature'') certifies optimality of one fixed-capacity inner
transportation LP, re-derived every outer iteration as capacities change; it
says nothing about the joint replica-and-routing strategy across the outer
loop. \textsc{FaaS-MAPG}'s certificate is a game-theoretic equilibrium
statement on exactly that joint strategy $(r_i,x_i,\omega_i,y_{i\cdot})_i$,
persisting across sweeps rather than being re-certified per inner solve, at
the cost of being restricted to the implemented move class rather than to
all deviations.

The exact-potential-game construction is standard once payoffs are chosen as
each player's contribution to a shared, additively separable objective
\citep{monderershapley1996}, and the underlying resource-sharing structure
(nodes competing for shared, capacitated neighbour capacity) is the setting
of congestion games \citep{rosenthal1973}, of which potential games are a
generalization admitting best/better-response convergence and the
finite-improvement property exploited by Theorem~\ref{thm:mapg-termination}.
The auction-based literature underlying \textsc{FaaS-MADeA}
\citep{bertsekas1988} and the dual-decomposition literature underlying
\textsc{FaaS-MALD} \citep{bertsekas1989,nedic2009} both target a different
certificate (price-based optimality of an inner allocation problem) rather
than a Nash-equilibrium notion on the full decentralized game; \textsc{FaaS-MAPG}
is, to our knowledge, the only family in this paper whose stopping condition
is an explicit game-theoretic equilibrium certificate.
